% =====================================================================
%  Lecture 2026-10-29
%
%  Reordered:
%   - the uniform boundedness principle is stated directly before its
%     proof, after Baire's theorem, which the proof uses; the source
%     states it before the category material and proves it after
%   - the product-space proposition is placed before the definition of
%     a closed operator, which uses the norm of X x Y
%   - the differentiation example, claimed before the product-space
%     proposition and verified after the sequential criterion, is given
%     once, after the criterion it uses
%
%  Normalized silently:
%   - B_r(x_0) written B(x_0, r), as everywhere else
%   - ||phi_N(g)|| and ||f(Tx) - f(T_n x)|| written with absolute values
%   - the question "What about R \ Q?" turned into a pointer to the
%     corollary that answers it
%   - the contradiction argument in the Fourier divergence proof put in
%     order
%   - the differentiation example on [0,1], where the verification
%     writes [a,b] with integrals from 0
%
%  CORRECTED SOURCE ERRORS
%   [p1] remark: "the other three do not have Banach in their name";
%        the uniform boundedness principle is also called the
%        Banach-Steinhaus theorem, and the open mapping theorem the
%        Banach-Schauder theorem  ->  only the correct content kept,
%        that Hahn-Banach needs no completeness and the others do
%   [p1] representation of (C[a,b])': w(t) = f~(1_[a,t]) for all t,
%        which shifts every Riemann-Stieltjes sum by g(a) w(a)  ->
%        w(a) = 0
%   [p1] geometric remark: "1 <= ||y|| for all y in V"  ->  for all
%        y in x_0 + V; ||x_0|| = 1 made explicit
%   [p1] Baire: every complete metric space  ->  every nonempty one,
%        the empty space being of first category
%   [p2] Dirichlet kernel: the underbrace labels 1/2 + sum cos(nt) as
%        D_N, while the closed form and the factor 1/(2 pi) need
%        D_N = 1 + 2 sum cos(nt)  ->  the latter
%   [p2] Fourier divergence: the proof runs in all of C[0,2 pi], while
%        the statement requires f(0) = f(2 pi)  ->  the proof runs in
%        the closed subspace of such functions
%   [p3] ONS example: <e_n, z> = z_n  ->  conj(z_n)
%   [p4] strong-limit proposition: hypothesis T_n -> T, which is
%        uniform convergence by the definition just given  ->  T_n x
%        convergent for every x, T the pointwise limit
%
%  ADDITIONS
%   [p1] representation of (C[a,b])': stated as a sharpening of
%        Theorem 11.10, which it repeats; BV[a,b] and B[a,b] named
%   [p1] interior glossed in the category definition
%   [p1] corollary: the union of two first-category sets is of first
%        category
%   [p2] UBP proof: continuity of T_n from Theorem 9.1; the case x = 0
%   [p2] Fourier coefficients defined, in the form of Example 7.8
%   [p3] weak convergence: Corollary 18.5 cited for uniqueness; Remark
%        18.4 cited in place of the source's repeated proof that j is
%        isometric
%   [p3] ONS example: the general case by Bessel's inequality, on
%        instruction
%   [p3] operator examples: strong convergence of T_n and ||T_n|| = 1
%        attained at e_{n+1}; ||S_n x|| = ||x|| for the failure of
%        strong convergence
%   [p4] strong-limit proposition: linearity of the limit, and the
%        passage ||Tx|| = lim ||T_n x||
%   [p4] bounded inverse: T^{-1} linear; the preimage under T^{-1} of U
%        is T(U); continuity gives boundedness by Theorem 9.1
%   [p4] sequential criterion: convergence in X x Y is convergence of
%        both components
%   [p4] differentiation: g continuous as a uniform limit (Theorem
%        4.13); limit and integral interchanged by uniform convergence
%
%  REPORTED BUT UNMODIFIED GAPS
%   [p1] representation of (C[a,b])': the proof is a sketch, as in the
%        source, and uses Hahn-Banach over K, proved here only over R
%   [p1] category examples asserted; Baire's theorem stated without
%        proof, on instruction
%   [p2] that sgn(D_N) can be approximated, and that the integral of
%        |D_N| tends to infinity, asserted, on instruction
%   [p4] open mapping theorem stated without proof, on instruction
%   [p4] product-space proposition: the completeness half is proved, by
%        the proof given in the lecture of 2026-11-05 (scan
%        funkanal_261105.pdf, p1, "Lemma 1"), moved here on instruction
%        so that statement and proof stay together; the norm axioms for
%        ||(x,y)|| = ||x|| + ||y|| are not checked
%
%  OMITTED
%   [p2] the small sketch of a ball inside A_{k_0}
%   [p4] the closed graph theorem and the first line of its proof, on
%        instruction; it is typeset with the next lecture
%   [p4] doodles
% =====================================================================

% [p1]
\begin{remark}
Of the four fundamental theorems of functional analysis, only the
Hahn--Banach theorem holds in every normed space; the other three require
Banach spaces.
\end{remark}

\begin{proposition}[Representation of {$(C[a,b])'$}, with norm]\label{prop:dual-C-norm}
Let $BV[a,b]$ denote the functions of bounded variation on $[a,b]$,
Definition~\ref{def:bv}. For every $f \in (\Cont{a,b})'$ the function
$w \in BV[a,b]$ of Theorem~\ref{thm:dual-C}, with
$f(g) = \int_a^b g(t) \dd w(t)$ for all $g \in \Cont{a,b}$, can be chosen so
that
\[
  \norm{f} = \operatorname{Var}(w)
  = \sup \sum_{j=1}^{N} \abs{w(t_j) - w(t_{j-1})} ,
\]
the supremum over all partitions $a = t_0 < t_1 < \dots < t_N = b$.
\end{proposition}

\begin{proof}[Proof (sketch)]
Suppose $f$ is given; we want to construct $w$. We have
\[
  (\Cont{a,b}, \norm{\cdot}_\infty) \subset (B[a,b], \norm{\cdot}_\infty),
  \qquad
  B[a,b] = \set{g \colon [a,b] \to \K \mid g \text{ bounded}} .
\]
By the
Hahn--Banach theorem, Theorem~\ref{thm:hb-normed}, there is an
$\tilde{f} \colon B[a,b] \to \K$ with $\tilde{f}|_{\Cont{a,b}} = f$ and
$\norm{f} = \norm{\tilde{f}}$. Then define $w(a) = 0$ and
$w(t) = \tilde{f}(\mathbf{1}_{[a,t]})$ for $t \in (a,b]$, where
$\mathbf{1}_{[a,t]}$ is the indicator function of $[a,t]$. Then
$w \colon [a,b] \to \K$, and it satisfies all of the requirements for $w$.
\end{proof}

\begin{remark}\label{rem:supporting}
Corollary~\ref{cor:norming-functional} states that for every $x_0 \ne 0$
there is an $f \in X'$ with $\norm{f} = 1$ and $f(x_0) = \norm{x_0}$. For a
geometric interpretation, consider the $p$-norms on $\R^d$. Let
$\norm{x_0} = 1$, let $f$ be the functional for $x_0$ from the corollary,
and let $V = \ker f$, a hyperplane. Every $y = x_0 + z$ with $z \in V$
satisfies $f(y) = f(x_0) + f(z) = \norm{x_0} = 1$. Also
$f(y) \le \norm{f}\,\norm{y} = \norm{y}$, so $1 \le \norm{y}$ for all
$y \in x_0 + V$: no point of the hyperplane $x_0 + V$ lies inside the unit
ball, and $x_0 + V$ touches the unit sphere at $x_0$.
\end{remark}

\begin{nfigure}{The hyperplane $x_0 + V$, $V = \ker f$, touching the unit
sphere of $\norm{\cdot}_4$ in $\R^2$ at $x_0$, with $\norm{x_0}_4 = 1$.}
\begin{tikzpicture}[scale=1.6]
  \draw[rulegrey, -{Stealth[length=1.6mm]}] (-1.4,0) -- (1.9,0);
  \draw[rulegrey, -{Stealth[length=1.6mm]}] (0,-1.4) -- (0,1.9);
  \draw[accent, thick, smooth, samples=120, domain=0:360]
    plot ({sign(cos(\x))*pow(abs(cos(\x)),0.5)},
          {sign(sin(\x))*pow(abs(sin(\x)),0.5)});
  \draw[thmaccent, thick] (0.18,1.50) -- (1.50,0.18);
  \draw[rulegrey, dashed] (-1.05,1.05) -- (1.05,-1.05);
  \fill (0.841,0.841) circle (0.025);
  \node[above right] at (0.85,0.85) {$x_0$};
  \fill (1.40,0.282) circle (0.025);
  \node[right] at (1.42,0.30) {$y$};
  \node[thmaccent] at (0.80,1.50) {$x_0 + V$};
  \node at (1.15,-1.05) {$V$};
\end{tikzpicture}
\end{nfigure}

\begin{nfigure}{The unit balls of $\norm{\cdot}_1$, $\norm{\cdot}_2$ and
$\norm{\cdot}_\infty$ in $\R^2$.}
\begin{tikzpicture}[scale=1.3]
  \draw[rulegrey, -{Stealth[length=1.6mm]}] (-1.4,0) -- (1.5,0);
  \draw[rulegrey, -{Stealth[length=1.6mm]}] (0,-1.4) -- (0,1.5);
  \draw[accent, thick] (1,0) -- (0,1) -- (-1,0) -- (0,-1) -- cycle;
  \draw[thmaccent, thick] (0,0) circle (1);
  \draw[exaccent, thick] (-1,-1) rectangle (1,1);
  \node[accent] at (0.33,0.33) {$1$};
  \node[thmaccent] at (0.62,0.62) {$2$};
  \node[exaccent] at (1.15,1.15) {$\infty$};
\end{tikzpicture}
\end{nfigure}

\begin{definition}[Category]\label{def:category}
In a metric space $(X,d)$ a subset $M \subset X$ is called
\begin{enumerate}
  \item \emph{nowhere dense}, if $\operatorname{int}(\cl{M}) = \emptyset$,
    where $\operatorname{int} A$, the interior of $A$, is the set of points
    of $A$ some ball around which lies in $A$;
  \item \emph{of first category}, if $M = \bigcup_{k=1}^{\infty} H_k$, where
    $H_k$ is nowhere dense for every $k$;
  \item \emph{of second category}, if it is not of first category.
\end{enumerate}
\end{definition}

\begin{example}
$\Z \subset \R$ is nowhere dense. $\Q \subset \R$ is not nowhere dense, but it
is of first category: $\Q = \bigcup_{i=1}^{\infty} \set{q_i}$. Whether
$\R \setminus \Q$ is of first category is settled by
Corollary~\ref{cor:irrationals}.
\end{example}

\begin{theorem}[Baire's category theorem]\label{thm:baire}
Every nonempty complete metric space $(X,d)$ is of second category in itself,
that is, if $X = \bigcup_{k=1}^{\infty} A_k$ where every $A_k$ is closed,
then $\operatorname{int} A_i \ne \emptyset$ for some $i$.
\end{theorem}

\begin{corollary}\label{cor:irrationals}
$\R \setminus \Q$ is of second category.
\end{corollary}

\begin{proof}
If it were of first category, we could write $\R = \Q \cup (\R \setminus \Q)$
as a union of two countable families of nowhere dense sets, which together
form one countable family, so $\R$ would be of first category. But $\R$
cannot be of first category, because of Theorem~\ref{thm:baire}.
\end{proof}

\begin{theorem}[Uniform boundedness principle]\label{thm:ubp}
Let $T_n \in \Bop{X}{Y}$, where $X$ is a Banach space and $Y$ is a normed
space, such that $(T_n)$ is \emph{pointwise bounded}: for every $x \in X$
there is a $c_x$ with $\norm{T_n x} \le c_x$ for all $n$. Then $(T_n)$ is
\emph{uniformly bounded}: there is a $c$ with $\norm{T_n} \le c$ for all
$n$.
\end{theorem}

% [p2]
\begin{proof}
Let $A_k = \set{x \in X \mid \norm{T_n x} \le k \ \forall n}$. Then $A_k$ is
closed: if $x_m \to x$ with $x_m \in A_k$, then $\norm{T_n x_m} \le k$ for
all $n$, and $T_n x_m \to T_n x$, $T_n$ being continuous by
Theorem~\ref{thm:bounded-continuous}, so $\norm{T_n x} \le k$. We cannot
converge out of $A_k$, so $A_k$ is closed.

$X = \bigcup_{k=1}^{\infty} A_k$ because $(T_n)$ is pointwise bounded: for
every $x \in X$ and any $k \ge c_x$ we have $x \in A_k$. $X$ is a Banach
space, so by Theorem~\ref{thm:baire} there is a $k_0$ with
$\operatorname{int} A_{k_0} \ne \emptyset$, hence
$B(x_0, r) \subset A_{k_0}$ for some $x_0$ and $r > 0$.

For $x = 0$ there is nothing to show. For $x \in X \setminus \set{0}$ let
$z = x_0 + \gamma x$ with $\gamma = \frac{r}{2\norm{x}}$. Then
$\norm{z - x_0} = \frac{r}{2}$, so $z \in B(x_0, r)$ and
$\norm{T_n z} \le k_0$ for all $n$; also $\norm{T_n x_0} \le k_0$. So
\[
  \norm{T_n x} = \Bigl\lVert T_n\Bigl(\frac{1}{\gamma}(z - x_0)\Bigr)
  \Bigr\rVert = \frac{1}{\gamma}\norm{T_n(z - x_0)}
  \le \frac{1}{\gamma}\bigl(\underbrace{\norm{T_n z}}_{\le k_0}
    + \underbrace{\norm{T_n x_0}}_{\le k_0}\bigr)
  \le \frac{2k_0}{\gamma} = \frac{4}{r} k_0 \norm{x} ,
\]
so $\norm{T_n} \le \frac{4}{r} k_0$ for all $n$.
\end{proof}

For a real $f \in \Cont{0,2\pi}$ let
\[
  a_0 = \frac{1}{2\pi}\int_0^{2\pi} f(t) \dd t , \qquad
  a_n = \frac{1}{\pi}\int_0^{2\pi} f(t) \cos(nt) \dd t , \qquad
  b_n = \frac{1}{\pi}\int_0^{2\pi} f(t) \sin(nt) \dd t , \qquad n \ge 1 ,
\]
so that $a_0 + \sum_{n=1}^{\infty} \bigl(a_n \cos(nt) + b_n \sin(nt)\bigr)$
is its Fourier series, as in Example~\ref{ex:fourier}.

\begin{proposition}\label{prop:fourier-divergence}
There is an $f \in \Cont{0,2\pi}$ with $f(0) = f(2\pi)$, but
\[
  f(t) \ne a_0 + \sum_{n=1}^{\infty} \bigl(a_n \cos(nt) + b_n \sin(nt)\bigr)
  \qquad \text{for } t = 0 ;
\]
that is, $f$ cannot be expanded in its Fourier series.
\end{proposition}

\begin{proof}
Let $X = \set{g \in \Cont{0,2\pi} \mid g \text{ real}, \ g(0) = g(2\pi)}$
with $\norm{\cdot}_\infty$. It is a closed subspace of a Banach space, by
Theorem~\ref{thm:C-sup-banach} with $[0,2\pi]$ in place of $[0,1]$, hence a
Banach space. Let $\varphi_N \colon X \to \R$ be the linear functional
defined, for $g \in X$ with Fourier coefficients $a_n, b_n$, by
\[
  \varphi_N(g) = \Bigl(a_0 + \sum_{n=1}^{N}
    \bigl(a_n \cos(nt) + b_n \sin(nt)\bigr)\Bigr)\Big|_{t=0}
  = a_0 + \dots + a_N
  = \frac{1}{2\pi}\int_0^{2\pi} g(t)\, D_N(t) \dd t ,
\]
where $D_N(t) = 1 + 2\sum_{n=1}^{N} \cos(nt)$ is the \emph{Dirichlet
kernel}, equal to $\sin((N + \frac{1}{2})t) / \sin(\frac{t}{2})$ for
$t \in (0, 2\pi)$. We show that $\varphi_N$ is bounded linear:
\[
  \abs{\varphi_N(g)}
  = \Bigl\lvert \frac{1}{2\pi}\int_0^{2\pi} g(t) D_N(t) \dd t \Bigr\rvert
  \le \frac{1}{2\pi}\int_0^{2\pi} \underbrace{\abs{g(t)}}_{\le
    \norm{g}_\infty} \abs{D_N(t)} \dd t
  \le \frac{\norm{g}_\infty}{2\pi}\int_0^{2\pi} \abs{D_N(t)} \dd t .
\]
So $\norm{\varphi_N} \le \frac{1}{2\pi}\int_0^{2\pi} \abs{D_N(t)} \dd t$,
with equality achieved with $g(t) = \operatorname{sgn}(D_N(t))$, and the
discontinuous $\operatorname{sgn}$ can be approximated, so that in fact
$\norm{\varphi_N} = \frac{1}{2\pi}\int_0^{2\pi} \abs{D_N(t)} \dd t$. But
\[
  \frac{1}{2\pi}\int_0^{2\pi} \abs{D_N(t)} \dd t
  \xrightarrow[N \to \infty]{} \infty .
\]
Assume, for a contradiction, that $g(0) = a_0 + \sum_{n=1}^{\infty} a_n$ for
every $g \in X$. In that case, for every $g \in X$ the sequence
$\varphi_N(g)$ converges, hence is bounded, so $(\varphi_N)$ is pointwise
bounded. But then by the uniform boundedness principle,
Theorem~\ref{thm:ubp}, there is a $c$ with $\norm{\varphi_N} \le c$ for all
$N$, a contradiction.
\end{proof}

% [p3]
\section{Types of convergence}

In a normed space $(X, \norm{\cdot})$, $x_n \to x$ means
$\norm{x_n - x} \to 0$. In $\Bop{X}{Y}$, $T_n \to T$ means
$\norm{T_n - T} \to 0$.

\begin{definition}[Weak convergence in $(X, \norm{\cdot})$]\label{def:weak}
A sequence $(x_n)$ \emph{converges to $x$ weakly}, written
$x_n \xrightarrow{w} x$, if $f(x_n) \to f(x)$ for every $f \in X'$.
\end{definition}

\begin{proposition}\label{prop:weak-limit}
If $x_n \xrightarrow{w} x$, then
\begin{enumerate}
  \item the weak limit is unique;
  \item $(\norm{x_n})$ is bounded.
\end{enumerate}
\end{proposition}

\begin{proof}
(1) Assume $x_n \xrightarrow{w} x$ and $x_n \xrightarrow{w} y$. Then
$f(x_n) \to f(x)$ and $f(x_n) \to f(y)$, therefore $f(x) = f(y)$ for every
$f \in X'$, and therefore $x = y$ by Corollary~\ref{cor:separating}.

(2) Consider the natural map $j \colon X \to X''$, $x \mapsto g_x$, with
$g_x(f) = f(x)$. By Remark~\ref{rem:j-isometric}, $\norm{g_x} = \norm{x}$:
$j$ is an isometric embedding. For $x_n$ consider $g_n = g_{x_n} \in X''$,
given by $g_n(f) = f(x_n)$. We know $\norm{g_n} = \norm{x_n}$, so seeing that
$(\norm{g_n})$ is bounded is sufficient. Now $g_n \in X'' = \Bop{X'}{\K}$, we
know that $X'$ is a Banach space by Corollary~\ref{cor:dual-banach}, and
$(g_n)$ is pointwise bounded: for every $f \in X'$, $g_n(f) = f(x_n)$ is
bounded, being convergent. By the uniform boundedness principle,
Theorem~\ref{thm:ubp}, $(\norm{g_n})$ is bounded.
\end{proof}

\begin{example}\label{ex:ons-weak}
Let $(e_n)$ be an orthonormal sequence in a Hilbert space $H$, for example
$e_n = (0, \dots, 0, 1, 0, \dots) \in \lp{2}$, with the $1$ in position $n$.
Then $e_n \xrightarrow{w} 0$. Indeed, for $f \in H'$ there is a $z \in H$
with $f = \ip{\cdot}{z}$, by Theorem~\ref{thm:riesz-hilbert}, and
$f(e_n) = \ip{e_n}{z} = \overline{\ip{z}{e_n}} \to 0$: by Bessel's
inequality, Corollary~\ref{cor:bessel}, the partial sums of
$\sum_n \abs{\ip{z}{e_n}}^2$ are bounded by $\norm{z}^2$, so the series
converges and its terms tend to $0$. In $\lp{2}$, with
$z = (z_1, z_2, \dots)$, this reads $\ip{e_n}{z} = \overline{z_n} \to 0$,
since $z \in \lp{2}$.
\end{example}

\begin{definition}[Types of operator convergence]\label{def:operator-convergence}
Let $T_n, T \in \Bop{X}{Y}$. We say that
\begin{enumerate}
  \item $T_n \to T$ \emph{uniformly}, if $\norm{T_n - T} \to 0$;
  \item $T_n \xrightarrow{s} T$ \emph{strongly}, if $T_n x \to Tx$ for every
    $x \in X$;
  \item $T_n \xrightarrow{w} T$ \emph{weakly}, if $f(T_n x) \to f(Tx)$ for
    every $x \in X$ and every $f \in Y'$.
\end{enumerate}
\end{definition}

\begin{proposition}\label{prop:convergence-implications}
$T_n \to T \implies T_n \xrightarrow{s} T \implies T_n \xrightarrow{w} T$.
\end{proposition}

\begin{proof}
For every $x$, $\norm{Tx - T_n x} \le \norm{T - T_n}\,\norm{x} \to 0$; and
for every $x$ and $f \in Y'$,
$\abs{f(Tx) - f(T_n x)} = \abs{f(Tx - T_n x)}
\le \norm{f}\,\norm{Tx - T_n x} \to 0$.
\end{proof}

\begin{example}\label{ex:tail-cut}
Let $T_n \colon \lp{2} \to \lp{2}$,
$(x_1, x_2, x_3, \dots) \mapsto (0, \dots, 0, x_{n+1}, x_{n+2}, \dots)$. Then
$T_n \xrightarrow{s} 0$, since $\norm{T_n x}^2 = \sum_{k > n} \abs{x_k}^2$ is
the tail of a convergent series. But $T_n \not\to 0$, because
$\norm{T_n} = 1$: $\norm{T_n x} \le \norm{x}$, and $T_n e_{n+1} = e_{n+1}$.
\end{example}

\begin{example}\label{ex:shift-weak}
Let $S_n \colon \lp{2} \to \lp{2}$,
$(x_1, x_2, \dots) \mapsto (0, \dots, 0, x_1, x_2, \dots)$, with $n$ zeros.
Then $S_n \xrightarrow{w} 0$, but $S_n \not\xrightarrow{s} 0$. For
$f \in (\lp{2})'$ write $f = \ip{\cdot}{z}$, by
Theorem~\ref{thm:riesz-hilbert}; then, with $T_n$ as in
Example~\ref{ex:tail-cut} and by the Cauchy--Schwarz inequality,
\[
  \abs{f(S_n x)} = \abs{\ip{S_n x}{z}}
  = \abs{x_1 \overline{z_{n+1}} + x_2 \overline{z_{n+2}} + \dots}
  \le \norm{x}\,\norm{T_n z} \longrightarrow 0 ,
\]
$\norm{T_n z}$ being the tail of a convergent series. On the other hand
$\norm{S_n x} = \norm{x}$, which does not tend to $0$ for $x \ne 0$.
\end{example}

% [p4]
\begin{proposition}\label{prop:strong-limit}
Let $T_n \in \Bop{X}{Y}$, where $X$ is a Banach space, and suppose that
$T_n x$ converges for every $x \in X$; let $Tx = \lim_{n \to \infty} T_n x$.
Then $T \in \Bop{X}{Y}$, and $T_n \xrightarrow{s} T$.
\end{proposition}

\begin{proof}
$T$ is linear, as a pointwise limit of linear maps. $(T_n)$ is pointwise
bounded: for every $x \in X$, $\norm{T_n x}$ is bounded because $T_n x$ is
convergent. By the uniform boundedness principle, Theorem~\ref{thm:ubp},
there is a $c$ with $\norm{T_n} \le c$. Therefore
$\norm{T_n x} \le c \norm{x}$ for all $n$, and letting $n \to \infty$,
$\norm{Tx} \le c\norm{x}$, so $\norm{T} \le c$.
\end{proof}

\section{Open mapping and closed graph theorems}

\begin{definition}[Open mapping]\label{def:open-mapping}
A mapping $T \colon (X,d) \to (Y,d')$ is an \emph{open mapping} if the image
$T(U)$ of every open set $U$ is also open.
\end{definition}

\begin{theorem}[Open mapping theorem]\label{thm:open-mapping}
Let $X$ and $Y$ be Banach spaces. Every $T \in \Bop{X}{Y}$ that is
surjective is an open mapping.
\end{theorem}

\begin{corollary}[Bounded inverse theorem]\label{cor:bounded-inverse}
If $T \in \Bop{X}{Y}$, $X$ and $Y$ are Banach spaces, and $T$ is bijective,
then $T^{-1} \in \Bop{Y}{X}$.
\end{corollary}

\begin{proof}
$X$ and $Y$ are Banach spaces and $T \in \Bop{X}{Y}$ is surjective, so by
Theorem~\ref{thm:open-mapping} $T$ is an open mapping. $T^{-1}$ is linear,
as the inverse of a linear bijection. For an open $U \subset X$, the preimage
of $U$ under $T^{-1}$ is $T(U)$, which is open. So $T^{-1}$ is continuous,
since the preimages of open sets are open, and therefore bounded, by
Theorem~\ref{thm:bounded-continuous}.
\end{proof}

\begin{proposition}\label{prop:product-space}
If $(X, \norm{\cdot})$ and $(Y, \norm{\cdot})$ are normed spaces, then
$(X \times Y, \norm{\cdot})$ is a normed space with
$\norm{(x,y)} = \norm{x} + \norm{y}$. If $X$ and $Y$ are Banach spaces, then
$X \times Y$ is also a Banach space.
\end{proposition}

% [scans/funkanal_261105.pdf, p1]
\begin{proof}
Let $X$ and $Y$ be Banach spaces, and take a Cauchy sequence
$(x_n, y_n) \subset X \times Y$. Then $(x_n)$ is Cauchy in $X$, since
$\norm{x_n - x_m} < \eps$ for all $n, m > N$ due to
\[
  \norm{(x_n - x_m, y_n - y_m)} = \norm{x_n - x_m} + \norm{y_n - y_m}
  < \eps .
\]
$(y_n)$ is also Cauchy, by the same argument. $X$ and $Y$ are Banach spaces,
so $x_n \to x$ and $y_n \to y$. Then $(x_n, y_n) \to (x, y)$.
\end{proof}

\begin{definition}[Closed operator]\label{def:closed-operator}
Let $X$ and $Y$ be normed spaces, and let $T \colon \operatorname{Dom} T \to
Y$, where $\operatorname{Dom} T \subset X$ is a subspace, and $T$ is linear.
$T$ is a \emph{closed operator} if its \emph{graph}
\[
  G(T) := \set{(x, Tx) \in X \times Y \mid x \in \operatorname{Dom} T}
\]
is closed in $X \times Y$, normed as in Proposition~\ref{prop:product-space}.
\end{definition}

\begin{proposition}\label{prop:closed-sequential}
Let $T \colon \operatorname{Dom} T \subset X \to Y$ be linear. Then $T$ is a
closed operator if and only if, for every sequence
$(x_n) \subset \operatorname{Dom} T$,
\[
  \bigl(x_n \to x \text{ and } T x_n \to y\bigr) \implies
  \bigl(x \in \operatorname{Dom} T \text{ and } Tx = y\bigr) .
\]
\end{proposition}

\begin{proof}
This is the ``converging out'' property: $G(T)$ is closed if and only if we
cannot converge out of it, meaning that if $(x_n, T x_n) \in G(T)$ and
$(x_n, T x_n) \to (x, y)$, then $(x, y) \in G(T)$, which just means
$x \in \operatorname{Dom} T$ and $y = Tx$. Convergence in $X \times Y$ is
convergence of both components, as
$\norm{(x_n, T x_n) - (x, y)} = \norm{x_n - x} + \norm{T x_n - y}$.
\end{proof}

\begin{example}[Differentiation]\label{ex:diff-closed}
Let $X = Y = \Cont{0,1}$ with $\norm{\cdot}_\infty$,
$\operatorname{Dom} T = \Contk{1}{0,1}$, and $Tf = f'$. We claim that $T$ is
a closed operator. By Proposition~\ref{prop:closed-sequential}, let
$f_n \in \Contk{1}{0,1}$ with $(f_n, f_n') \to (f, g)$, that is,
$\norm{f - f_n}_\infty \to 0$ and $\norm{g - f_n'}_\infty \to 0$; we need
$f \in \Contk{1}{0,1}$ and $f' = g$. The function $g$ is continuous, as a
uniform limit of continuous functions, Theorem~\ref{thm:C-sup-banach}. For
$t \in [0,1]$, interchanging limit and integral by uniform convergence,
\[
  \int_0^t g(\tau) \dd\tau = \int_0^t \lim_{n \to \infty} f_n'(\tau) \dd\tau
  = \lim_{n \to \infty} \int_0^t f_n'(\tau) \dd\tau
  = \lim_{n \to \infty} \bigl(f_n(t) - f_n(0)\bigr) = f(t) - f(0) .
\]
So $f(t) = f(0) + \int_0^t g(\tau) \dd\tau$, hence $f \in \Contk{1}{0,1}$ and
$f'(t) = g(t)$.
\end{example}
